\documentclass[11pt]{article}
\usepackage[T1]{fontenc}
\usepackage[margin=1in]{geometry}
\usepackage{amsmath,amssymb,amsthm,booktabs,microtype}
\usepackage[hidelinks]{hyperref}
\newtheorem{theorem}{Theorem}[section]
\newtheorem{proposition}[theorem]{Proposition}
\newtheorem{lemma}[theorem]{Lemma}
\newtheorem{corollary}[theorem]{Corollary}
\theoremstyle{definition}
\newtheorem{remark}[theorem]{Remark}
\newcommand{\Q}{\mathbb Q}
\newcommand{\Z}{\mathbb Z}
\newcommand{\C}{\mathbb C}
\newcommand{\Sym}{\operatorname{Sym}}
\newcommand{\diag}{\operatorname{diag}}
\newcommand{\Res}{\operatorname{res}}
\newcommand{\BR}{\mathrm{BR}}
\numberwithin{equation}{section}
\title{The Broadhurst--Roberts Recursion and Bessel Period Pairings}
\author{Alexander Kolpakov\\[0.5ex]
\small Center for Science, Technology, Engineering, and Mathematics\\
\small University of Austin, Austin, Texas, United States\\
\small \texttt{akolpakov@uaustin.org}\\[0.75ex]
\small The Wolfram Institute for Computational Foundations of Science\\
\small \texttt{akolpakov@wolframinstitute.org}}
\date{Research draft, 20 October 2026}
\begin{document}
\maketitle
\begin{abstract}
Broadhurst and Roberts gave a recursion for rational matrices appearing
in quadratic relations between Bessel moments. Fres\'an, Sabbah and Yu
proved the period relations with a geometric intersection matrix,
leaving its identification with that recursion open.
We identify the recursion matrix with coefficients of the
symmetric-power Bessel operator and then with the inverse de Rham
pairing. The argument includes the resonant symmetric powers.
In particular, the inverse middle pairing multiplied by $N!$ has
integer entries, answering a second question of Fres\'an--Sabbah--Yu.
The main step sums the auxiliary recursion in a period-ratio coordinate;
formal residues then recover the intersection form.
\end{abstract}

\section{Introduction and statements}
Put $\theta_z=z\,d/dz$ and consider the rank-two connection
\begin{equation}\label{eq:connection}
 \theta_z v_0=v_1,\qquad \theta_z v_1=zv_0.
\end{equation}
This is the Kloosterman connection $\mathrm{Kl}_2$ in the normalization
of Fres\'an--Sabbah--Yu \cite{FSY}. For $N\ge3$, set
\[
 r=\left\lfloor\frac{N-1}{2}\right\rfloor,\qquad q=r+1.
\]
Let $S_N$ denote their middle de Rham intersection matrix in the basis
of \cite[Corollary 3.11]{FSY}. If $4\nmid N$, its basis is
$\omega_i=[z^iv_0^N\,dz/z]$, $1\le i\le r$.
If $N=4\ell$, remove index $\ell$ and use
\begin{equation}\label{eq:middlebasis}
 \omega'_i=\omega_i-\gamma_{N,i}\omega_\ell,\qquad
 i\in I_N=\{1,\ldots,r\}\setminus\{\ell\},
\end{equation}
where
\begin{equation}\label{eq:gamma}
 2^N\bigl(I_0(2\sqrt z)K_0(2\sqrt z)\bigr)^{N/2}
 \sim\sum_{j\ge\ell}\gamma_{N,j}z^{-j},\qquad
 \gamma_{N,\ell}=1.
\end{equation}
For $4\nmid N$, let $I_N=\{1,\ldots,r\}$.
Indices in a submatrix always retain their original order.
For $N=4$ the middle matrix is empty; the identities below use the
usual zero-dimensional convention.

Define differential operators in $\Z[z,\theta_z]$ by
\begin{equation}\label{eq:operatorrecursion}
 P^-_0=1,\quad P^-_1=\theta_z,\qquad
 P^-_{k+1}=\theta_zP^-_k-k(N+1-k)zP^-_{k-1},\quad 1\le k\le N.
\end{equation}
Write $\mathcal L_N=P^-_{N+1}$ and form the integer matrix
\begin{equation}\label{eq:E}
 E_N(a,b)=[z^a]\mathcal L_Nz^{-b},\qquad 1\le a,b\le r.
\end{equation}
The operator recursion is established: it is the
Bronstein--Mulders--Weil construction, as stated and proved in
Borwein--Salvy \cite[Lemmas 2--3]{BS}.

\begin{theorem}\label{thm:pairing}
With the above normalization,
\begin{equation}\label{eq:mainpairing}
 N!S_N^{-1}=(-1)^NE_N[I_N,I_N].
\end{equation}
In particular $N!S_N^{-1}$ has integer entries for every $N\ge3$.
\end{theorem}

Broadhurst--Roberts \cite[formulas (5.7)--(5.12)]{BR} define a rational
$r\times r$ matrix $D_N^{\BR}$ through a bivariate recursion, recalled
in Section~\ref{sec:series}. Put
$T_N=\diag((-4)^1,\ldots,(-4)^r)$ and $T_I=T_N[I_N,I_N]$.

\begin{theorem}\label{thm:BR}
The literal Broadhurst--Roberts recursion satisfies
\begin{align}
 D_N^{\BR}&=T_N^{-1}E_NT_N^{-1},\label{eq:BRfull}\\
 D_N^{\BR}[I_N,I_N]
   &=(-1)^NN!\,(T_IS_NT_I)^{-1}.\label{eq:BRmiddle}
\end{align}
\end{theorem}

Equation~\eqref{eq:BRmiddle} proves the matrix identifications left
open in \cite[Remark 5.12]{FSY}. Equation~\eqref{eq:mainpairing}
answers the integrality question in the same remark.
The quadratic period relations themselves were proved in \cite{FSY},
with a different proof in Zhou \cite{Zhou}. Zhou also distinguishes
the recursion question in \cite[Section 4.3]{Zhou}.
Our argument supplies that identification, using their established
geometric pairing and residue formulas. It does not address the
conjecture that quadratic relations exhaust all algebraic relations
between these periods.

\section{Summing the auxiliary recursion}\label{sec:series}

\subsection{The local period coordinate}
Set $x=t^2/4$ and
\[
 J(x)=J_0(2\sqrt x)=\sum_{n\ge0}\frac{(-1)^nx^n}{(n!)^2},\qquad
 \theta=x\frac d{dx}.
\]
Thus $(\theta^2+x)J=0$. Define formal series
\begin{equation}\label{eq:CVW}
 C(x)=\frac{2(1-J^2+2xJJ')}{3x^2},\qquad
 V=\frac{J^2}{C}=\sum_{n\ge0}\frac{v_nx^n}{n!},\qquad
 W=-\frac{JJ'}{C}=\sum_{n\ge0}\frac{w_nx^n}{n!}.
\end{equation}
They have constant term one. These are exactly the series in
\cite[formulas (5.7)--(5.9)]{BR}: the change of variable uses
$J'_0(t)=-J_1(t)$.

There is a unique $\beta\in x\Q[[x]]$ such that
\begin{equation}\label{eq:beta}
 \beta'=\frac{J^{-2}-1}{x}.
\end{equation}
Put $w=\log x+\beta$. The two solutions $J,Jw$ have
$\theta$-Wronskian one, since $\theta w=J^{-2}$.
In particular,
\begin{equation}\label{eq:coordinate}
 Q=x\exp\beta,\qquad
 \frac{dQ}{Q}=\frac{dx}{xJ^2},\qquad
 Q\frac d{dQ}=J^2\theta.
\end{equation}
Here $Q=x+O(x^2)$ has a formal compositional inverse.

\subsection{The bivariate polynomials}
The polynomials of \cite[formula (5.10)]{BR} have $H_0=1$, $H_{-1}=0$,
and, for $s\ge1$,
\begin{equation}\label{eq:H}
 H_s(Y,Z)=ZH_{s-1}-(s-1)YH_{s-2}
 -\sum_{k=1}^{s-1}\binom{s-1}{k}
       \bigl(v_kH_{s-k}-w_kZH_{s-k-1}\bigr).
\end{equation}
Set
\begin{equation}\label{eq:dD}
 d_s(N,c)=\frac{H_s(3c/2,N+2-2c)}{4^ss!},\qquad
 D_N^{\BR}(a,b)=\sum_{c=-b}^{a}
     d_{a-c}(N,-c)d_{b+c}(N,c)c^{N+1}.
\end{equation}

\begin{lemma}\label{lem:sum}
For every integer $c$,
\begin{equation}\label{eq:sum}
 \sum_{s\ge0}d_s(N,c)(4x)^s
     =J^{-N-2}\exp(-c\beta).
\end{equation}
\end{lemma}
\begin{proof}
Let $H(x)=\sum_sH_s(Y,Z)x^s/s!$.
Multiply \eqref{eq:H} by $x^{s-1}/(s-1)!$ and sum.
The convolution terms are $(V-1)H'$ and $Z(W-1)H$, giving
\[
 VH'=(ZW-Yx)H.
\]
From \eqref{eq:CVW} and \eqref{eq:beta},
\[
 \frac WV=-\frac{J'}J,\qquad
 \frac{xC}{J^2}=\frac23\left(\beta'+2\frac{J'}J\right).
\]
For $Y=3c/2$ and $Z=N+2-2c$, it follows that
$H'/H=-(N+2)J'/J-c\beta'$. The initial value is $H(0)=1$,
which proves \eqref{eq:sum}.
\end{proof}

\begin{proposition}\label{prop:coefficient}
Define
\begin{equation}\label{eq:Bol}
 L_N=J^{-N-2}(J^2\theta)^{N+1}J^{-N}.
\end{equation}
Then, for all $a,b\ge1$,
\begin{equation}\label{eq:coefficient}
 D_N^{\BR}(a,b)=4^{-a-b}[x^a]L_Nx^{-b}.
\end{equation}
\end{proposition}
\begin{proof}
Put $f=J^{-N-2}$. Lemma~\ref{lem:sum} rewrites $4^{a+b}D_N^{\BR}(a,b)$
as the coefficient of $x^a$ in
\[
 f(x)\sum_{c\in\Z}c^{N+1}Q(x)^cg_c,\qquad
 g_c=[y^{b+c}]f(y)\exp(-c\beta(y)).
\]
Only $-b\le c\le a$ contributes. In the residue defining $g_c$,
make the formal substitution $Q=y\exp\beta(y)$.
Equation~\eqref{eq:coordinate} gives
\[
 g_c=\Res_{y=0}f(y)y^{-b}Q^{-c}\frac{dy}{y}
    =[Q^c]y(Q)^{-b}J(y(Q))^{-N}.
\]
Multiplication of the $c$th coefficient by $c^{N+1}$ is the operator
$(Q\,d/dQ)^{N+1}$. Substituting back and using
\eqref{eq:coordinate} gives exactly \eqref{eq:Bol}.
All operations concern formal Laurent coefficients; no summability
assumption is involved.
\end{proof}

\section{The symmetric-power operator}

\begin{lemma}\label{lem:operator}
The operator $L_N$ is the monic symmetric power of $\theta^2+x$.
It belongs to $\Z[x,\theta]$, has $x$-degree at most $q$, and
\begin{equation}\label{eq:adjoint}
 L_N^*=(-1)^{N+1}L_N
\end{equation}
for the formal adjoint associated with $dx/x$.
Moreover $\mathcal L_N=L_N|_{x=-z}$.
\end{lemma}
\begin{proof}
For completeness, apply the established recursion to
$u_k=y^{N-k}(\theta y)^k$, where $(\theta^2+x)y=0$.
Since $\theta u_k=(N-k)u_{k+1}-kx u_{k-1}$, the operators
\[
 P_0=1,\quad P_1=\theta,\qquad
 P_{k+1}=\theta P_k+k(N+1-k)xP_{k-1}
\]
satisfy $P_ku_0=N!u_k/(N-k)!$ for $k\le N$ and $P_{N+1}u_0=0$.
Varying $y$ through linear combinations of $J$ and $Jw$ shows that
$P_{N+1}$ kills their full symmetric power.
The same is true of \eqref{eq:Bol}: multiplication by $J^{-N}$
gives a polynomial of degree at most $N$ in $w$, and $J^2\theta w=1$.
Both operators are monic of order $N+1$. Their common solution space
has dimension $N+1$, so they coincide.
This proves the integer-coefficient assertion and the substitution
claim. Induction in the recursion bounds the weight of a term
$x^j\theta^k$ by $2j+k\le N+1$.

To compute the adjoint, put $A=J^2\theta$. Since $\theta^*=-\theta$,
$A^*=-J^{-2}AJ^2$. Reverse the factors in \eqref{eq:Bol} and cancel
the intervening powers of $J$. This gives \eqref{eq:adjoint}.
\end{proof}

The monic operator is independent of the chosen solution frame:
every change of frame preserves the symmetric-power solution space.
Thus the expression in local periods \eqref{eq:Bol} descends to the
integer operator \eqref{eq:operatorrecursion}.
The local series in the recursion determine coefficients of this
common invariant. This is the useful period-frame step of the proof.

\begin{lemma}\label{lem:leading}
Write $\mathcal L_N=\sum_{j=0}^qz^jp_j(\theta_z)$. Then
\begin{equation}\label{eq:leading}
 p_q(\theta_z)=
 \begin{cases}
 (-1)^q((2q-1)!!)^2,&N=2q-1,\\
 (-1)^q(2^qq!)^2(\theta_z+q/2),&N=2q.
 \end{cases}
\end{equation}
The matrix $E_N$ satisfies
\[
 E_N^T=(-1)^{N+1}E_N,\qquad E_N(a,b)=0\quad(a+b>q).
\]
It is invertible if $4\nmid N$. If $N=4\ell$, it has rank $r-1$,
and its principal submatrix on $I_N$ is invertible.
\end{lemma}
\begin{proof}
For $N=2q-1$, the top coefficient of $z^j$ in $P^-_{2j}$ is
$(-1)^j\prod_{h=1}^j(2h-1)(N+2-2h)$, directly from the recursion.
This gives the first formula.

For $N=2q$, the weight bound gives $\deg p_q\le1$.
In the associated commutative symbol, derivative commutators have
lower weight. The recurrence becomes the determinant recurrence for
the tridiagonal symmetric-power matrix with off-diagonal entries
$N-a$ and $az$. Its eigenvalues are $(N-2a)\sqrt z$,
$0\le a\le N$: these are the eigenvalues of the symmetric power of
the matrix with eigenvalues $\sqrt z,-\sqrt z$.
The coefficient of $z^q\theta_z$ in its characteristic polynomial is
$(-1)^q(2^qq!)^2$. Adjointness requires
$p_q(-\theta_z-q)=-p_q(\theta_z)$, fixing its constant term.

The adjoint identity, applied to the constant term of
$z^{-a}\mathcal L_Nz^{-b}$, proves the stated symmetry of $E_N$.
The degree bound gives its zero entries. On the anti-diagonal
$a+b=q$ its entries are $p_q(-b)$.
They are nonzero for odd $N$ and for $N\equiv2\pmod4$.
When $N=4\ell$, only the central anti-diagonal entry vanishes;
deleting its index leaves an invertible anti-triangular matrix.
The full matrix is skew-symmetric of odd size, so its rank is
exactly $r-1$.
\end{proof}

Proposition~\ref{prop:coefficient} and the substitution $x=-z$ now give
\begin{equation}\label{eq:ED}
 E_N(a,b)=(-4)^{a+b}D_N^{\BR}(a,b).
\end{equation}
This proves \eqref{eq:BRfull}.

\section{Formal primitives and the intersection matrix}

We use the existence of formal primitives and the residue description
of the pairing from \cite[Lemma 3.3, Proposition 3.8 and Section 3.b]{FSY}.
We recall the normalization to keep track of the sign and transpose.
For $u_a=v_0^{N-a}v_1^a$, their algebraic pairing is
\begin{equation}\label{eq:algebraicpairing}
 \langle u_a,u_b\rangle=
 \begin{cases}(-1)^a/\binom Na,&a+b=N,\\0,&a+b\ne N.
 \end{cases}
\end{equation}
The primitive at zero has coefficients in $z\C[[z]]$ for the
positive-power sources used here, so its residue contribution is zero.

Suppose a formal primitive at infinity is
$m=\sum_{a=0}^Nf_au_a$ and
$\nabla m=g(z)u_0\,dz/z$. Its components satisfy
\begin{align}\label{eq:components}
 \theta_zf_a+(N-a+1)f_{a-1}+(a+1)zf_{a+1}&=0
                      &&(1\le a\le N),\\
 \theta_zf_0+zf_1&=g.\nonumber
\end{align}
Put $\eta=f_N$.

\begin{lemma}\label{lem:primitive}
These equations imply
\begin{equation}\label{eq:elimination}
 \mathcal L_N\eta=(-1)^NN!g.
\end{equation}
If $\deg g\le q-1$, every formal Laurent solution $\eta$ at infinity
has only strictly negative powers of $z$.
\end{lemma}
\begin{proof}
The equations with $a\ge1$ determine the other components from $\eta$.
The coefficient of $\theta_z^{N+1}\eta$ in the last equation is
$(-1)^N/N!$. The homogeneous eliminant has as solutions the top
components of horizontal symmetric-power sections.
For a horizontal rank-two section, its top component solves
$(\theta_z^2-z)y=0$. Products of those sections therefore give exactly
the full symmetric-power solution space in the top component.
The monic eliminant is consequently $\mathcal L_N$, proving
\eqref{eq:elimination} including its factor.

Write the highest exponent of a nonzero Laurent series $\eta$ as $M$.
If $M\ge0$, Lemma~\ref{lem:leading} gives $p_q(M)\ne0$.
The highest exponent on the left of \eqref{eq:elimination} is then
$M+q>\deg g$, a contradiction.
\end{proof}

Let $M_{bi}=[z^{-b}]\eta_i$ for a family of primitives with sources
$g_i$. A term $z^{-b}$ with $b\ge q$ cannot contribute to a positive
power under $\mathcal L_N$. The lemma therefore gives the finite identity
\begin{equation}\label{eq:finite}
 E_NM=(-1)^NN!G,
\end{equation}
where $G_{ai}=[z^a]g_i$, $1\le a\le r$.

\subsection{Nonresonant powers}
Suppose $4\nmid N$ and take $g_i=z^i$.
Then $G=1$ and $M=(-1)^NN!E_N^{-1}$.
By \eqref{eq:algebraicpairing}, the residue at infinity is
\begin{equation}\label{eq:residue}
 S_N(i,j)=(-1)^{N+1}M_{ji};
\end{equation}
this is \cite[formula (3.15)]{FSY}.
Thus
\[
 S_N=-N!E_N^{-T},\qquad N!S_N^{-1}=(-1)^NE_N,
\]
where the second equality uses Lemma~\ref{lem:leading}.

\subsection{Resonant powers}
Suppose $N=4\ell$ and let
$h=(\gamma_{N,1},\ldots,\gamma_{N,r})^T$.
The formal series in \eqref{eq:gamma} solves $\mathcal L_Nh(z)=0$:
it is the product of $N/2$ copies of each scalar Bessel solution.
Taking positive coefficients gives $E_Nh=0$.
Since $h_\ell=1$, Lemma~\ref{lem:leading} shows that this vector
spans the kernel.

Let $U$ be the $r\times(r-1)$ matrix whose columns are
\begin{equation}\label{eq:U}
 U_i=e_i-h_ie_\ell,\qquad i\in I_N.
\end{equation}
For the middle sources $g_i=z^i-h_iz^\ell$,
equation~\eqref{eq:finite} becomes $E_NM=N!U$.
The ambiguity of a primitive is a multiple of the horizontal formal
section whose top component is \eqref{eq:gamma}; see
\cite[Proposition 3.8]{FSY}. We may therefore impose the gauge
$M_{\ell i}=0$.
Writing $\bar E=E_N[I_N,I_N]$ gives
\[
 M[I_N,I_N]=N!\bar E^{-1}.
\]
The residue with the corrected source is
\[
 S_N(i,j)=-(M_{ji}-h_jM_{\ell i})=-M_{ji}.
\]
Adding a horizontal section does not change it, since
$h_j-h_jh_\ell=0$. Thus the gauge computes the same pairing as
\cite[formula (3.16)]{FSY}, independently of their normalization of
the primitive. We obtain
\[
 S_N=-N!\bar E^{-T}=N!\bar E^{-1},\qquad N!S_N^{-1}=\bar E.
\]
This proves Theorem~\ref{thm:pairing} in all cases.
Combining it with \eqref{eq:ED} proves \eqref{eq:BRmiddle} and
Theorem~\ref{thm:BR}.

\section{The quadratic relation, including the singular matrix}

For clarity we also recover the original normalization.
Let $M(a,b,c)=\int_0^\infty I_0(t)^aK_0(t)^bt^c\,dt$ and
$Z(s)=|B_s|/s!$, with $B_s$ the Bernoulli number.
The matrices in \cite[formulas (5.2)--(5.6)]{BR} are
\begin{align*}
 F_{2r+1}(u,a)&=\frac{(-1)^{a-1}}{\pi^u}
                  M(r+1-u,r+u,2a-1),\\
 F_{2r+2}(u,a)&=\frac{(-1)^{a-1}}{\pi^{u+1/2}}
                  M(r+1-u,r+1+u,2a-1),\\
 B_{2r+1}^{\BR}(u,v)&=(-1)^{u+r}2^{-2r-2}
                  (r+u)!(r+v)!Z(u+v),\\
 B_{2r+2}^{\BR}(u,v)&=(-1)^{u+r}2^{-2r-3}
                  (r+u+1)!(r+v+1)!Z(u+v+1),
\end{align*}
with all indices from $1$ to $r$.

\begin{corollary}\label{cor:quadratic}
For every $N\ge3$, with the literal recursion matrix,
\[
 F_ND_N^{\BR}F_N^T=B_N^{\BR}.
\]
\end{corollary}
\begin{proof}
If $4\nmid N$, this follows at once from Theorem~\ref{thm:BR} and
\cite[Corollary 5.10]{FSY}.
We give the pullback argument needed for $N=4\ell$.
Write $P$ and $B$ for the $r\times r$ matrices of
\cite[Sections 4--5]{FSY}, retaining indices $1,\ldots,r$.
The period matrix on the middle sources, with all these homology rows,
is $Q=PU$. Both its rows and the pairing $B$ factor through the map
from rapid-decay homology to middle homology:
$\alpha_i\mapsto\beta_i$, as in \cite[formula (4.5)]{FSY}.
The basis of the latter space is $\beta_2,\ldots,\beta_r$.
If $A$ records the coordinates of all $\beta_i$ in that basis, then
\[
 Q=AP^{\rm mid},\qquad B=AB^{\rm mid}A^T.
\]
The established middle period relation
\cite[Theorem 5.3]{FSY} consequently gives
\[
 QS_N^{-1}Q^T=(-2\pi\mathrm i)^{N+1}B.
\]
The identity $E_Nh=0$, with $h_\ell=1$, implies
$E_N=U\bar EU^T$: the blocks on $I_N\times I_N$ agree, and the
kernel determines the remaining row and column. Since
$S_N^{-1}=\bar E/N!$, we obtain
\[
 PE_NP^T=N!(-2\pi\mathrm i)^{N+1}B.
\]
Finally the normalizations in \cite[formula (5.9)]{FSY} are
\[
 F_N=(-2\sqrt\pi)^{-N-1}R^TPT_N,\qquad
 B_N^{\BR}=-\frac{\mathrm i^{N+1}N!}{2^{N+1}}R^TBR,
\]
where $R=\operatorname{antidiag}(\mathrm i,\ldots,\mathrm i^r)$.
Substitute \eqref{eq:BRfull}. The scalar factor is
$(-2\sqrt\pi)^{-2N-2}N!(-2\pi\mathrm i)^{N+1}
=-\mathrm i^{N+1}N!/2^{N+1}$ because $N$ is even.
This proves the full singular relation as well.
\end{proof}

\section{Examples and the role of the period frame}

For $N=5,6$, the formulas give
\[
 D_5^{\BR}=\begin{pmatrix}13/8&225/64\\225/64&0\end{pmatrix},\quad
 S_5=\begin{pmatrix}0&8/15\\8/15&208/3375\end{pmatrix},
\]
\[
 D_6^{\BR}=\begin{pmatrix}0&-18\\18&0\end{pmatrix},\qquad
 S_6=\begin{pmatrix}0&-5/8\\5/8&0\end{pmatrix}.
\]
The pairings reproduce \cite[Examples 3.19--3.20]{FSY}.
For $N=8$, the first nontrivial resonant example is
\[
 D_8^{\BR}=\begin{pmatrix}0&-288&-576\\288&0&0\\576&0&0\end{pmatrix},
 \qquad S_8=\begin{pmatrix}0&35/128\\-35/128&0\end{pmatrix},
\]
with $I_8=\{1,3\}$ and $h=(0,1,1/8)^T$.
The omitted row and column reflect a horizontal formal solution,
not a failure of the coefficient formula.

The proof separates a local series presentation from the invariant
it determines. A determinant-one frame $(J,Jw)$ belongs to the
torsor of normalized solution frames of the rank-two connection.
The period ratio supplies $Q$, in which the convolution in
\eqref{eq:dD} becomes an Euler derivative.
Returning to $x$ gives \eqref{eq:Bol}; frame invariance then identifies
it with the integer symmetric-power operator. The established
residue pairing turns its finite coefficients into $S_N^{-1}$.
Thus the frame comparison connects the local recursion with a global
intersection form, identifying the Broadhurst--Roberts recursion
and proving the factorial-scaled integrality in
Theorems~\ref{thm:pairing}--\ref{thm:BR}.

This is a separate application of the period-frame approach used for
quartic periods in \cite{Quartic}; here it is applied to the
symmetric-power Bessel connection.

The accompanying exact checker compares the literal recursion with
\eqref{eq:operatorrecursion} for $3\le N\le32$.
An independent downward elimination in the full connection checks
the residue signs and normalization, including the corrected sources
and kernel in every resonant case in that range.
All factorial-scaled inverse matrices are integral, as the theorem
requires. These finite checks supplement the formal proofs above.
The source and data are in \path{research/bessel-recursion};
\texttt{make bessel-recursion-check} reproduces them.

\bibliographystyle{abbrvurl}
\bibliography{bessel-recursion}
\end{document}
